\scanpage{99}

\deleted{Si $N$ est \unsure{équivalente à la}} \deleted{Sous ces conditions,}

Alors la fonction $u\to N(u)$ sur $L^N_{\deleted{\unsure{0}}}(E,F)$ est une \uline{norme}. En effet, on a $N(u+v)=\lim_{n\to\infty}N(|(\rho_i(u+v))|_n)\deleted{\unsure{,\rho_2(u+v)}}$, or le second mb est majoré par $N(|(\rho_i(u))|_n)\deleted{\unsure{,\rho_2(u)}}+N(|(\rho_i(v))|_n)\deleted{\unsure{,\rho_2(v)}}$ (N\textsuperscript{o} , formule (44)), donc par $N(u)+N(v)$. On note aussi trivialement $N(\lambda u)=|\lambda|N(u)$, et $N(u)=0\Rightarrow u=0$.

\margin{$N(u)=N(u^*)$ (12)\par $\left.\begin{aligned}N(uv)&\leq\lVert u\rVert N(v)\\N(vu)&\leq\lVert u\rVert N(v)\end{aligned}\right\}$ (13)}

L'espace $L^N_{\deleted{\unsure{0}}}(E,F)$ est donc un espace normé; \deleted{\unsure{prouvons qu'il est \uline{complet}.}}\inserted{soit} $L_N(E,F)$ le sous-espace normé formé des $u$ tels que $\unsure{(\rho_i(u))|(\rho_i u)}\in\ell_N$. Je dis que $L_N(E,F)$ est l'adhérence de $\overline E\otimes F$ dans $L^N(E,F)$. En effet, montrons que

$\sum\rho_i\bar e_i\otimes f_i$ est limite au sens de $N$ de $\sum_1^n\rho_i\bar e_i\otimes f_i$, donc que $N(\sum_{n+1}^\infty\rho_i\bar e_i\otimes f_i)=N(\rho_{n+1}\unsure{(u)},\ldots)\to 0$ pour $n\to+\infty$, en d'autres termes, que dans $\ell_N$ les opérations $x\to t_n(x)=(\overbrace{0{-}0}^{n},x_{n+1},x_{n+2}\unsure{-})$ tendent vers $0$ simplement. Or elles forment une suite équicontinue, qui tend vers zéro dans \deleted{\unsure{un}} $R^{(N)}$, ce qui prouve la chose. \deleted{\unsure{Notons} \unsure{donc|$N_\infty$}, \unsure{lorsqu'on a} $L^N(E,F)$ \unsure{et si|et que} $N$ \illegible{3}}

\margin{Définition de $L^N(E,F)$ si $N\simeq N_\infty$ ($L^N(E,F)=L(E,F)$, avec \textcircled{14} $\boxed{N(u)=\underset{\substack{v\in F'\otimes E\\N^\circ(v)\leq 1}}{\operatorname{Sup}}\operatorname{Tr}(vu)}$ $L_N(E,F)$ comme avant). Cas où $u$ est cpt.}

\deleted{Le dual de $L_N(E,F)$ s'identifie à un sous-espace du dual de $L_{N_1}(E,F)=L^1(E,F)$, donc de $L(F,E)$, supposons que $N$ n'est pas équivalent à $N_1$, donc $N'$ non équivalent à $N_\infty$. Je dis que}

\[
\deleted{(L_N(E,F))'=L^{N'}(F,E).}
\]

\deleted{En premier lieu, \deleted{\unsure{si}} si $u\in L_N(E,F)$, $v\in L^{N'}(F,G)$ alors $uv\in L^1(E,G)$, avec}

\[
\deleted{\lVert uv\rVert_1\leq N(u)N'(v)}
\]

\deleted{soit} Soient $N,N'\ N''$ des normes de Schatten telles que

\[
\boxed{N(xy)\leq N'(x)N''(y)}\tag*{(\unsure{15|13})}
\]

pour \unsure{tout|tt} $x,y\in R^{(N)}$. Alors la $=$ formule sera valable pour $x\deleted{\unsure{,y|y}}\in\ell^{N'}$, $y\in\ell^{N''}$. Prouvons alors que $u\in L^{N'}(E,F)$, $v\in L^{N''}(F,G)$ implique $uv\in L^N(E,G)$, avec

\[
\boxed{N(uv)\leq N'(u)N''(v)}\tag*{(16)}
\]

\margin{100}
